\documentclass[11pt,a4paper,oneside]{article}

% --- Поддержка русского языка в pdfLaTeX ---
\usepackage[T2A]{fontenc}
\usepackage[utf8]{inputenc}
\usepackage[russian]{babel}

% --- Математические библиотеки ---
\usepackage{amsmath,amssymb,amsfonts,amsthm,mathtools}

% --- Геометрия страницы и типографика ---
\usepackage[left=20mm, right=20mm, top=22mm, bottom=22mm, headheight=14pt]{geometry}
\usepackage{microtype}
\usepackage{tabularx}
\usepackage{array}
\usepackage{booktabs}
\usepackage{verbatim}
\usepackage{xcolor}
\usepackage{fancyhdr}
\usepackage{hyperref}

\hypersetup{
    colorlinks=true,
    linkcolor=blue!70!black,
    urlcolor=blue!70!black,
    citecolor=blue!70!black,
    pdftitle={Задача B. Калибровка магического альтиметра},
    pdfauthor={Яндекс CodeRun}
}



% --- Настройка олимпиадных макросов Polygon / Codeforces ---
\newenvironment{problem}[5]{
    \addcontentsline{toc}{section}{#1}
    \begin{center}
        {\LARGE\bfseries #1}\par
        \vspace{0.3cm}
        {\small\sffamily
        \begin{tabular}{rl}
            \textbf{Входной файл:} & #2 \\
            \textbf{Выходной файл:} & #3 \\
            \textbf{Ограничение по времени:} & #4 \\
            \textbf{Ограничение по памяти:} & #5 \\
        \end{tabular}
        }\par
    \end{center}
    \vspace{0.3cm}
    \par\noindent
}{
    \vspace{0.8cm}
}

% --- Секции олимпиадного условия ---
\newcommand{\InputFile}{\subsubsection*{Формат входных данных}}
\newcommand{\OutputFile}{\subsubsection*{Формат выходных данных}}
\newcommand{\Examples}{\subsubsection*{Примеры}}
\newcommand{\Example}{\subsubsection*{Пример}}
\newcommand{\Note}{\subsubsection*{Примечания}}

% --- Колонтитулы ---
\pagestyle{fancy}
\fancyhf{}
\fancyhead[L]{\small\sffamily Задача B. Калибровка магического альтиметра}
\fancyhead[R]{\small\sffamily Яндекс CodeRun (2025-winter-common)}
\fancyfoot[C]{\thepage}
\renewcommand{\headrulewidth}{0.4pt}
\renewcommand{\footrulewidth}{0pt}

\setlength{\parindent}{1.5em}
\setlength{\parskip}{0.4ex plus 0.2ex minus 0.1ex}

\begin{document}

% ----------------------------------------------------------------------
% Задача B. Калибровка магического альтиметра
% Тема: Динамическое программирование и рюкзак | День: 2
% ----------------------------------------------------------------------
\begin{problem}{B. Калибровка магического альтиметра}{стандартный ввод}{стандартный вывод}{1.0 секунда}{512 МБ}

Кодерун хочет поскорее взойти на~Пик Кода и~спасти цифровой мир от~Вечной зимы. Чтобы покорение высоты прошло без проблем, сперва нужно откалибровать магический альтиметр~---~устройство, измеряющее высоту над~уровнем цифрового моря. Альтиметр нужен при~ориентировании в~горах, в~условиях плохой видимости, он оповещает о~перепадах высот и~о достижении заданной точки.

На~приборе имеется 10 шкал, на~каждой из~которых выставлено целое число от~1 до~100. Нужно выбрать подмножество этих шкал (возможно, пустое), чтобы сумма выбранных чисел была как можно ближе к~100.

Если несколько сумм находятся на~одинаковом расстоянии от~100, следует выбрать наибольшую из~них. Например, для сумм 98 и~102 нужно вывести 102.

\InputFile
Вводится 10 чисел от~1 до~100, каждое в~отдельной строке~---~значения на~шкалах альтиметра.

\OutputFile
Выведите ответ на~задачу~---~сумму чисел на~выбранных шкалах~---~результат калибровки магического альтиметра.

\Examples
\begin{center}
\begin{tabular}{|p{0.48\textwidth}|p{0.48\textwidth}|}
\hline
\textbf{Стандартный ввод} & \textbf{Стандартный вывод} \\
\hline
\begin{minipage}[t]{0.47\textwidth}
\vspace{0pt}
\begin{verbatim}
10
20
30
40
50
60
70
80
90
100
\end{verbatim}
\vspace{0pt}
\end{minipage}
&
\begin{minipage}[t]{0.47\textwidth}
\vspace{0pt}
\begin{verbatim}
100
\end{verbatim}
\vspace{0pt}
\end{minipage} \\
\hline
\multicolumn{2}{|p{0.96\textwidth}|}{\small\textit{Пояснение к примеру 1:} Шкала со значением 100 дает точную сумму 100.} \\
\hline
\begin{minipage}[t]{0.47\textwidth}
\vspace{0pt}
\begin{verbatim}
1
2
3
5
8
13
21
34
55
89
\end{verbatim}
\vspace{0pt}
\end{minipage}
&
\begin{minipage}[t]{0.47\textwidth}
\vspace{0pt}
\begin{verbatim}
100
\end{verbatim}
\vspace{0pt}
\end{minipage} \\
\hline
\multicolumn{2}{|p{0.96\textwidth}|}{\small\textit{Пояснение к примеру 2:} Подмножество выбранных чисел дает в~сумме ровно 100.} \\
\hline
\begin{minipage}[t]{0.47\textwidth}
\vspace{0pt}
\begin{verbatim}
40
40
40
40
40
40
40
40
40
40
\end{verbatim}
\vspace{0pt}
\end{minipage}
&
\begin{minipage}[t]{0.47\textwidth}
\vspace{0pt}
\begin{verbatim}
120
\end{verbatim}
\vspace{0pt}
\end{minipage} \\
\hline
\multicolumn{2}{|p{0.96\textwidth}|}{\small\textit{Пояснение к примеру 3:} Возможные суммы кратны 40. Ближайшие к~100 суммы~---~это 80 (разность 20) и~120 (разность 20). По~правилу выбора наибольшей выводится 120.} \\
\hline
\end{tabular}
\end{center}

\Note
Пустое множество шкал дает сумму 0.

\end{problem}


\end{document}
